\documentclass[../../main.tex]{subfiles}% 注意这里的文件路径不能用 ./main.tex ,否则用latexmk编译子文件会报错
\graphicspath{{\subfix{./image/}}{\subfix{../../image/}}} % 指定图片引用路径,后续可以直接使用图片文件名
% 如果图片存放在上级目录,则使用 ../../image/ 作为备用路径

% 例如:
% \begin{figure}[H]
% \centering
% \includegraphics[width=0.3\textwidth]{图.png}
% \caption{}
% \label{figure:图}
% \end{figure}
% 注意:上述\label{}一定要放在\caption{}之后,否则引用图片序号会只会显示??.

\begin{document}

\section{最大模原理和Schwarz引理}

\begin{theorem}[最大模原理]\label{theorem:最大模原理-定理4.5.1}
设 \( f \) 是区域 \( D \) 中非常数的全纯函数，那么 \( |f(z)| \) 不可能在 \( D \)的内部取到最大值。
\end{theorem}
\begin{proof}

{\color{blue}证法一:}因为 \( f \) 是 \( D \) 上非常数的全纯函数，由\refthe{theorem:定理4.4.6-开映射定理}，\( G = f(D) \) 是一个区域。如果 \( |f(z)| \) 在 \( D \) 中某点 \( z_0 \) 处取到最大值，记 \( w_0 = f(z_0)\in G \)，则由$G$是区域知 \( w_0 \) 是 \( G \) 的一个内点，即有 \( \varepsilon > 0 \)，使得 \( B(w_0, \varepsilon) \subseteq G \)。故必有 \( w_1 \in G \)，使得 \( |w_1| > |w_0| \)。于是存在 \( z_1 \in D \)，使得 \( |f(z_1)| = |w_1| > |w_0| = |f(z_0)| \)。这与 \( |f(z_0)| \) 是 \( |f(z)| \) 在 \( D \) 中的最大值相矛盾！

{\color{blue}证法二:}假设存在 \(z_0 \in D\) 使得 \(|f(z_0)| = M = \max_{z \in D} |f(z)|\)。由于 \(D\) 是区域，存在 \(r > 0\) 使得闭圆盘 \(\overline{B(z_0, r)} \subseteq D\)。对任意 \(\rho \in (0, r]\)，函数 \(f\) 在圆 \(|\zeta - z_0| < \rho\) 内解析，在闭圆 \(|\zeta - z_0| \leqslant \rho\) 上连续，故由\hyperref[theorem:解析函数的平均值定理]{解析函数的平均值定理}可得
\begin{align*}
f(z_0) = \frac{1}{2\pi} \int_{0}^{2\pi} f(z_0 + \rho \mathrm{e}^{\mathrm{i}\varphi}) \mathrm{d}\varphi.
\end{align*}
取模并利用三角不等式，得
\begin{align}\label{eqLLjjifiowjeiofjew}
M = |f(z_0)| \leqslant \frac{1}{2\pi} \int_{0}^{2\pi} |f(z_0 + \rho \mathrm{e}^{\mathrm{i}\varphi})| \mathrm{d}\varphi \leqslant M\Longrightarrow \frac{1}{2\pi} \int_{0}^{2\pi} |f(z_0 + \rho \mathrm{e}^{\mathrm{i}\varphi})| \mathrm{d}\varphi=M.
\end{align}
由于 \(|f(z_0 + \rho \mathrm{e}^{\mathrm{i}\varphi})| \leqslant M\)对一切 \(\varphi \in [0, 2\pi]\)成立.若存在$\varphi_0\in [0,2\pi]$使得$|f(z_0 + \rho \mathrm{e}^{\mathrm{i}\varphi_0})| < M$.不妨设$\varphi_0\in (0,2\pi)$,则由$|f(z_0 + \rho \mathrm{e}^{\mathrm{i}\varphi})|$关于$\varphi$连续知,存在$\delta>0$,使得
\begin{align*}
|f(z_0 + \rho \mathrm{e}^{\mathrm{i}\varphi})|<M,\quad \forall \varphi\in (\varphi_0-\delta,\varphi_0+\delta).
\end{align*}
从而
\begin{align*}
\frac{1}{2\pi}\int_0^{2\pi}{\left| f\left( z_0+\rho \mathrm{e}^{\mathrm{i}\varphi} \right) \right|\mathrm{d}\varphi}&=\frac{1}{2\pi}\left( \int_0^{\varphi _0-\delta}{\left| f\left( z_0+\rho \mathrm{e}^{\mathrm{i}\varphi} \right) \right|\mathrm{d}\varphi}+\int_{\varphi _0-\delta}^{\varphi _0+\delta}{\left| f\left( z_0+\rho \mathrm{e}^{\mathrm{i}\varphi} \right) \right|\mathrm{d}\varphi}+\int_{\varphi _0+\delta}^{2\pi}{\left| f\left( z_0+\rho \mathrm{e}^{\mathrm{i}\varphi} \right) \right|\mathrm{d}\varphi} \right) 
\\
&<\frac{1}{2\pi}\left( 2\delta M+M\left( 2\pi -2\delta \right) \right) =M,
\end{align*}
这与\eqref{eqLLjjifiowjeiofjew}式矛盾!
于是对一切 \(\varphi \in [0, 2\pi]\) 必有
\begin{align*}
|f(z_0 + \rho \mathrm{e}^{\mathrm{i}\varphi})| = M.
\end{align*}
故对任意 \(\rho \in (0, r]\)，圆周 \(|z - z_0| = \rho\) 上均有 \(|f(z)| = M\)。从而由$\rho$的任意性知在圆盘 \(B(z_0, r)\) 内，\(|f(z)|\) 恒等于 \(M\)。由\cref{proposition:复变函数导数为0必是常函数}，故 \(f\) 在 \(B(z_0, r)\) 内为常数。再由\hyperref[theorem:复变函数--唯一性定理]{唯一性定理}，\(f\) 在区域 \(D\) 内为常数，这与 \(f\) 非常数矛盾！因此 \(|f(z)|\) 不可能在 \(D\) 的内部取到最大值。
\end{proof}

\begin{theorem}\label{theorem:定理4.5.2}
设 \( D \) 是 \( \mathbb{C} \) 中的有界区域，如果非常数的函数 \( f \) 在 \( \overline{D} \) 上连续，在 \( D \) 内全纯，那么 \( f \) 的最大模在 \( D \) 的边界上而且只在 \( D \) 的边界上达到。
\end{theorem}
\begin{remark}
\refthe{theorem:定理4.5.2}中 \( D \) 的有界性条件不能去掉，否则定理可能不成立。例如，设
\[
D = \left\{ z : |\text{Im} z| < \frac{\pi}{2} \right\},
\quad
f(z) = \mathrm{e}^{\mathrm{e}^z}.
\]
当然 \( f \) 在 \( D \) 中全纯，在 \( \overline{D} \) 上连续，但它的最大模并不能在 \( \partial D \) 上达到。事实上，当 \( z \in \partial D \) 时，\( z = x \pm \frac{\pi}{2}\mathrm{i} \)，这时，\( \mathrm{e}^z = \mathrm{e}^x \mathrm{e}^{\pm \frac{\pi}{2}\mathrm{i}} = \pm \mathrm{i}\mathrm{e}^x \)，所以 \( |\mathrm{e}^{\mathrm{e}^z}| = |\mathrm{e}^{\pm \mathrm{i}\mathrm{e}^x}| = 1 \)。而当 \( z \in D \) 时，取 \( z = x \)，即有 \( \mathrm{e}^{\mathrm{e}^x} \to \infty \)（\( x \to \infty \)）。故\refthe{theorem:定理4.5.2}不成立。
\end{remark}
\begin{proof}
因为 \( \overline{D} \) 是紧集，其上的连续函数 \( |f| \) 一定有最大值，即存在 \( z_0 \in \overline{D} \)，使得 \( |f(z_0)| \) 是 \( |f(z)| \) 在 \( \overline{D} \) 上的最大值。由\refthe{theorem:最大模原理-定理4.5.1} 知道，\( z_0 \) 不能属于 \( D \)，因此只能有 \( z_0 \in \partial D \)。 
\end{proof}

\begin{theorem}[Schwarz引理]\label{theorem:Schwarz引理-定理4.5.3}
设 \( f \) 是单位圆盘 \( B(0,1) \) 中的全纯函数，且满足条件

(i) 当 \( z \in B(0,1) \) 时，\( |f(z)| \leqslant  1 \)；

(ii) \( f(0) = 0 \)，

那么下列结论成立：

(i) 对于任意 \( z \in B(0,1) \)，均有 \( |f(z)| \leqslant  |z| \)；

(ii) \( |f'(0)| \leqslant  1 \)；

(iii) 存在某点 \( z_0 \in B(0,1) \)，\( z_0 \neq 0 \)，使得 \( |f(z_0)| = |z_0| \)，或者 \( |f'(0)| = 1 \) 成立的充要条件是存在实数 \( \theta \)，使得对 \( B(0,1) \) 中所有的 \( z \)，都有 \( f(z) = \mathrm{e}^{\mathrm{i}\theta} z \)。
\end{theorem}
\begin{proof}
因为 \( f \in H(B(0,1)) \)，且 \( f(0) = 0 \)，故 \( f \) 在 \( B(0,1) \) 中可展开为
\begin{align*}
f(z) = a_1 z + a_2 z^2 + \cdots = z(a_1 + a_2 z + \cdots) = z g(z),
\end{align*}
这里，\( g(0) = a_1 = f'(0) \)。取 \( 0 < r < 1 \)，当 \( |z| = r \) 时，有
\[
|g(z)| = \frac{|f(z)|}{|z|} \leqslant  \frac{1}{r},
\]
故由\hyperref[theorem:最大模原理-定理4.5.1]{最大模原理}，在圆盘 \( B(0,r) \) 中也有
\[
|g(z)| \leqslant  \frac{1}{r} \, (\text{当} \ |z| < r \ \text{时}).
\]
让 \( r \to 1 \)，即得 \( |g(z)| \leqslant  1 \)（\( z \in B(0,1) \)），即 \( |f(z)| \leqslant  |z| \)，结论(i)成立。

从 \( |g(0)| \leqslant  1 \) 即得 \( |f'(0)| \leqslant  1 \)，结论(ii)成立。

结论(iii)的充分性显然,只证必要性.现若有 \( z_0 \in B(0,1) \)，\( z_0 \neq 0 \)，使得 \( |f(z_0)| = |z_0| \)，即 \( |g(z_0)| = 1 \)。这说明全纯函数 \( g \) 在内点 \( z_0 \) 处取到了最大模 1，根据\hyperref[theorem:最大模原理-定理4.5.1]{最大模原理}，\( g \) 必须是常数。设 \( g(z) \equiv c \)，由 \( |g(z_0)| = 1 \)，得 \( |c| = 1 \)，所以 \( c = \mathrm{e}^{\mathrm{i}\theta} \)，因而 \( f(z) = \mathrm{e}^{\mathrm{i}\theta} z \)。如果 \( |f'(0)| = 1 \)，即 \( |g(0)| = 1 \)，与上面一样讨论，即得 \( f(z) = \mathrm{e}^{\mathrm{i}\theta} z \)。结论(iii)成立。
\end{proof}

\begin{corollary}[更精确的Schwarz引理]\label{theorem:更精确的Schwarz引理}
设 \( f \) 是单位圆盘 \(\mathbb{D} = B(0,1) \) 中的全纯函数，且满足条件

(i) 当 \( z \in B(0,1) \) 时，\( |f(z)| \leqslant  1 \)；

(ii) \( f(0) = 0 \)，

则
\begin{align}\label{eq::2-jfo8sj8fj238fjowj30wf}
\left| f\left( z \right) \right|\leqslant \left| z \right|\frac{\left| z \right|+\left| f' \left( 0 \right) \right|}{1+\left| f'\left( 0 \right) \right|\left| z \right|},\quad \forall z\in \mathbb{D}.
\end{align}
\end{corollary}
\begin{proof}
令
\begin{align*}
g(z)=
\begin{cases}
\dfrac{f(z)}{z},\quad z\neq 0,\\
f'(0),\quad z=0
\end{cases}
\end{align*}
则$g\in H(\mathbb{D})$。由\hyperref[theorem:Schwarz引理-定理4.5.3]{Schwarz引理}知
\begin{align*}
|g(0)|=|f'(0)|\leqslant 1,\quad |g(z)|=\frac{|f(z)|}{|z|}\leqslant 1,\quad \forall z\in\mathbb{D}\setminus\{0\}.
\end{align*}
若存在$z_0\in\mathbb{D}$，使得$|g(z_0)|=1$，则由\hyperref[theorem:Schwarz引理-定理4.5.3]{Schwarz引理}知存在$\theta\in\mathbb{R}$，使得$f(z)=e^{\mathrm{i}\theta}z$，此时$|f'(0)|=1$，代入\eqref{eq::2-jfo8sj8fj238fjowj30wf}式知，此时\eqref{eq::2-jfo8sj8fj238fjowj30wf}式成立。下设$|g(z)|<1,\forall z\in\mathbb{D}$。不妨设$g(0)=a=|f'(0)|\in[0,1)$，否则用$e^{-\mathrm{i}\arg f'(0)}g(z)$代替$g(z)$。

再令$\varphi(z)=\frac{a-z}{1-az}$，则$\varphi\in\mathrm{Aut}(\mathbb{D}),\varphi(a)=0$。从而
\begin{align*}
\varphi\circ g(0)=0,\quad |\varphi\circ g(z)|\leqslant 1,\quad \forall z\in\mathbb{D}.
\end{align*}
由\hyperref[theorem:Schwarz引理-定理4.5.3]{Schwarz引理}知，对$\forall z\in\mathbb{D}$，有
\begin{align*}
|\varphi\circ g(z)|\leqslant |z|\Longleftrightarrow \left|\frac{a-g(z)}{1-ag(z)}\right|\leqslant |z|.
\end{align*}
从而对$\forall z\in\mathbb{D}$，有
\begin{align*}
\left|\frac{a-g(z)}{1-ag(z)}\right|^2\leqslant |z|^2\Longleftrightarrow (1-a^2|z|^2)|g(z)|^2+a^2-|z|^2\leqslant 2a(1-|z|^2)\mathrm{Re}g(z).
\end{align*}
再由$\mathrm{Re}g(z)\leqslant |g(z)|$可得
\begin{align*}
(1-a^2|z|^2)|g(z)|^2-2a(1-|z|^2)|g(z)|+(a^2-|z|^2)\leqslant 0\Longleftrightarrow |g(z)|\in\left[\frac{a-|z|}{1-a|z|},\frac{a+|z|}{1+a|z|}\right].
\end{align*}
因此
\begin{align*}
|g(z)|\leqslant \frac{a+|z|}{1+a|z|}\Longleftrightarrow |f(z)|\leqslant |z|\frac{|z|+|f'(0)|}{1+|f'(0)||z|},\quad \forall z\in\mathbb{D}.
\end{align*}
\end{proof}

\begin{theorem}[Osserman定理]\label{theorem:Osserman定理-2000年}
设\(\mathbb{D} = B(0,1) \),$f:\mathbb{D}\to\mathbb{D}$ 全纯, $f(0)=0$. 若 $f$ 可连续至边界点 $b\in\partial\mathbb{D}$, 且 $|f(b)|=1$, $f'(b)$ 存在, 则
\begin{align*}
|f'(b)|\geqslant \frac{2}{1+|f'(0)|}.
\end{align*}
特别地, $|f'(b)|\geqslant 1$, 等号成立当且仅当存在实数$\theta$,使得$f(z)=e^{\mathrm{i}\theta}z$.
\end{theorem}
\begin{proof}
由\hyperref[theorem:更精确的Schwarz引理]{更精确的Schwarz引理}知
\begin{align*}
|f(z)|\leqslant |z|\frac{|z|+|f'(0)|}{1+|f'(0)||z|},\quad \forall z\in\mathbb{D}.
\end{align*}
再结合条件$|b|=1,|f(b)|=1$可得
\begin{align*}
\left|\frac{f(z)-f(b)}{|z|-|b|}\right|\geqslant \frac{1-|f(z)|}{1-|z|}\geqslant \frac{1-|z|\frac{|z|+|f'(0)|}{1+|f'(0)||z|}}{1-|z|}=\frac{1+|z|}{1+|f'(0)||z|},\quad \forall z\in\mathbb{D}.
\end{align*}
由于$f'(b)$存在，故上式两边令$z$沿径向趋于$b$，得到
\begin{align}\label{eq::-290jr0wjf0j34jf0jew0fejw90}
|f'(b)|=\lim\limits_{z\rightarrow b}\left|\frac{f(z)-f(b)}{|z|-|b|}\right|\geqslant \frac{2}{1+|f'(0)|}.
\end{align}
特别地,由\hyperref[theorem:Schwarz引理-定理4.5.3]{Schwarz引理}知$|f'(0)|\leqslant 1$,等号成立当且仅当存在实数$\theta$,使得$f(z)=e^{\mathrm{i}\theta}z$.再由\eqref{eq::-290jr0wjf0j34jf0jew0fejw90}式可得
\begin{align*}
|f'(b)|\geqslant \frac{2}{1+|f'(0)|}\geqslant 1.
\end{align*}
此时等号成立也当且仅当存在实数$\theta$,使得$f(z)=e^{\mathrm{i}\theta}z$.
\end{proof}

\begin{definition}
设 \( D \) 是 \( \mathbb{C} \) 中的区域，如果 \( f \) 是 \( D \) 上的单叶全纯函数，且 \( f(D) = D \)，就称 \( f \) 是 \( D \) 上的一个\textbf{全纯自同构}。\( D \) 上全纯自同构的全体记为 \( \text{Aut}(D) \)。
\end{definition}

\begin{proposition}
\( \text{Aut}(D) \) 在复合运算下构成一个群，称为 \( D \) 的\textbf{全纯自同构群}。
\end{proposition}
\begin{proof}
设 \( f, g \in \text{Aut}(D) \)，那么 \( f \circ g \in \text{Aut}(D) \)，且复合运算满足结合律。对于每个 \( f \in \text{Aut}(D) \)，由\refthe{theorem:定理4.4.9}，\( f^{-1} \in \text{Aut}(D) \)。\( f(z) = z \) 在复合运算下起着单位元素的作用。因而 \( \text{Aut}(D) \) 在复合运算下构成一个群.
\end{proof}

对于一般的区域 \( D \)，要确定 \( \text{Aut}(D) \) 是很困难的。但对于单位圆盘 \( B(0,1) \)，应用 \hyperref[theorem:Schwarz引理-定理4.5.3]{Schwarz引理}不难定出其上的全纯自同构群。

对于 \( |a| < 1 \)，记
\[
\varphi_a(z) = \frac{a - z}{1 - \overline{a} z},
\]
由\refpro{proposition:例2.5.16} 知道，它把 \( B(0,1) \) 一一地映为 \( B(0,1) \)，因而 \( \varphi_a \in \text{Aut}(B(0,1)) \)。如果记 \( \rho_\theta(z) = \mathrm{e}^{\mathrm{i}\theta} z \)，它是一个旋转变换，当然有 \( \rho_\theta \in \text{Aut}(B(0,1)) \)。下面我们将证明，\( \text{Aut}(B(0,1)) \) 中除了 \( \varphi_a, \rho_\theta \) 以及它们的复合外，不再有其他的变换。

\begin{theorem}\label{theorem:定理4.5.5}
设 \( f \in \text{Aut}(B(0,1)) \)，且 \( f^{-1}(0) = a \)，则必存在 \( \theta \in \mathbb{R} \)，使得
\[
f(z) = \mathrm{e}^{\mathrm{i}\theta} \frac{a - z}{1 - \overline{a} z}.
\]
\end{theorem}
\begin{proof}
记 \( w = \varphi_a(z)= \frac{a - z}{1 - \overline{a} z} \)，直接计算可得
\begin{align}\label{equation::::14865}
z = \varphi_a^{-1}(w) = \frac{a - w}{1 - \overline{a} w} = \varphi_a(w).
\end{align}
令 \( g(w) = f \circ \varphi_a(w) \)，则由\refpro{proposition:例2.5.16}知道 \( g \in \text{Aut}(B(0,1)) \)，而且
\[
g(0) = f(\varphi_a(0)) = f(a) = 0,
\]
故由 \hyperref[theorem:Schwarz引理-定理4.5.3]{Schwarz引理}得
\begin{align}
|g'(0)| \leqslant  1.\label{eq::::1--1---1-1-1-1}
\end{align}
由于 \( g^{-1} \in \text{Aut}(B(0,1)) \)，且 \( g^{-1}(0) = 0 \)，故对 \( g^{-1} \) 用\hyperref[theorem:Schwarz引理-定理4.5.3]{Schwarz引理}，得 \( |(g^{-1})'(0)| \leqslant  1 \)。但由\refthe{theorem:定理4.4.9}，有
\[
|(g^{-1})'(0)| = \frac{1}{|g'(0)|},
\]
由此即得
\[
|g'(0)| \geqslant  1.
\]
与 \(\eqref{eq::::1--1---1-1-1-1}\) 式比较，即得 \( |g'(0)| = 1 \)。根据 \hyperref[theorem:Schwarz引理-定理4.5.3]{Schwarz引理的结论(iii)}，存在实数 \( \theta \)，使得 \( g(w) = \mathrm{e}^{\mathrm{i}\theta} w \)，即 \( f \circ \varphi_a(w) = \mathrm{e}^{\mathrm{i}\theta} w \)。令 \( w = \varphi_a(z) \)，再结合\eqref{equation::::14865}式即得
\[
f(z) = \mathrm{e}^{\mathrm{i}\theta} \frac{a - z}{1 - \overline{a} z}.
\] 
\end{proof}

\begin{theorem}[Schwarz-Pick定理]\label{theorem:Schwarz-Pick定理-定理4.5.6}
设 \( f: B(0,1) \to B(0,1) \) 是全纯函数，对于 \( a \in B(0,1) \)，\( f(a) = b \)。那么

(i) 对任意 \( z \in B(0,1) \)，有 \( |\varphi_b(f(z))| \leqslant  |\varphi_a(z)| \),也即
\begin{align*}
\left|\frac{f(a) - f(z)}{1 - \overline{f(a)}f(z)}\right| \leqslant \left|\frac{a - z}{1 - \overline{a} z}\right|,
\end{align*}
其中$\varphi_a(z) = \frac{a - z}{1 - \overline{a} z},\varphi_b(z) = \frac{b - z}{1 - \overline{b} z}$;

(ii) \( |f'(a)| \leqslant  \frac{1 - |b|^2}{1 - |a|^2} \);

(iii) 如果存在某点 \( z_0 \in B(0,1) \)，\( z_0 \neq a \)，使得 \( |\varphi_b(f(z_0))| = |\varphi_a(z_0)| \)，或者 \( |f'(a)| = \frac{1 - |b|^2}{1 - |a|^2} \) 成立，
那么 \( f \in \text{Aut}(B(0,1)) \)。其中$\varphi_a(z) = \frac{a - z}{1 - \overline{a} z},\varphi_b(z) = \frac{b - z}{1 - \overline{b} z}.$
\end{theorem}
\begin{proof}
令 \( g = \varphi_b \circ f \circ \varphi_a \)，则 \( g \in H(B(0,1)) \)，且 \( g(B(0,1)) \subset B(0,1) \)，\( g(0) = \varphi_b \circ f \circ \varphi_a(0) = 0 \)。对 \( g \) 用\hyperref[theorem:Schwarz引理-定理4.5.3]{Schwarz引理}，有
\begin{align}
|\varphi_b \circ f \circ \varphi_a(\zeta)| \leqslant  |\zeta|, \, \zeta \in B(0,1) \label{eq::::--2-2-2-2-2}
\end{align}
和
\begin{align}
|(\varphi_b \circ f \circ \varphi_a)'(0)| \leqslant  1. \label{eq::::--2-2-2-2-3}
\end{align}
令 \( z = \varphi_a(\zeta) \)，则 \( \zeta = \varphi_a(z) \)，于是 \(\eqref{eq::::--2-2-2-2-2}\) 式变成
\begin{align}
|\varphi_b(f(z))| \leqslant  |\varphi_a(z)|. \label{eq::::--2-2-2-2-4}
\end{align}
这就是(i)。

由于
\[
\varphi_a'(0) = - (1 - |a|^2),
\quad
\varphi_b'(b) = - \frac{1}{1 - |b|^2},
\]
由 \(\eqref{eq::::--2-2-2-2-3}\) 式即得
\begin{align}
|f'(a)| \leqslant  \frac{1 - |b|^2}{1 - |a|^2}. \label{eq::::--2-2-2-2-5}
\end{align}
这就是(ii)。

如果存在 \( z_0 \in B(0,1) \)，\( z_0 \neq a \)，使得 \(\eqref{eq::::--2-2-2-2-4}\) 式中的等号成立，令 \( \zeta_0 = \varphi_a(z_0) \)，则 \( \zeta_0 \neq 0 \)，且 \( \zeta_0 \) 使 \(\eqref{eq::::--2-2-2-2-2}\) 式中的等号成立。于是由 \hyperref[theorem:Schwarz引理-定理4.5.3]{Schwarz引理}，\( g(z) = \mathrm{e}^{\mathrm{i}\theta} z \)，即 \( g \in \text{Aut}(B(0,1)) \)，于是 \( f = \varphi_b \circ g \circ \varphi_a \in \text{Aut}(B(0,1)) \)。

注意到当\(\eqref{eq::::--2-2-2-2-5}\) 式中的等号成立时,有
\begin{align*}
g' (0)&=\varphi _{b}'\left[ f\left( \varphi _a\left( 0 \right) \right) \right] \cdot f'\left( \varphi _a\left( 0 \right) \right) \cdot \varphi _{a}'\left( 0 \right) 
=\varphi _{b}'\left[ f\left( a \right) \right] \cdot f' \left( a \right) \cdot \varphi _{a}'\left( 0 \right) 
\\
&=\varphi _{b}'\left[ b \right] \cdot f' \left( a \right) \cdot \varphi _{a}'\left( 0 \right) 
=-\frac{1}{1-|b|^2}\cdot \frac{1-|b|^2}{1-|a|^2}\cdot \left[ -\left( 1-|a|^2 \right) \right] =1,
\end{align*}
由 \hyperref[theorem:Schwarz引理-定理4.5.3]{Schwarz引理}，\( g(z) = \mathrm{e}^{\mathrm{i}\theta} z \)，即 \( g \in \text{Aut}(B(0,1)) \)，于是 \( f = \varphi_b \circ g \circ \varphi_a \in \text{Aut}(B(0,1)) \)。
\end{proof}

\begin{example}
设$\mathbb{D} =\left\{ z:\left| z \right|<1 \right\}$，$f:\mathbb{D} \rightarrow \mathbb{D}$全纯，非恒等映射.
\begin{enumerate}[(1)]
\item $f$至多有一个不动点;
\item 举例说明$f$可能无不动点.
\end{enumerate}
\end{example}
\begin{proof}
\begin{enumerate}[(1)]
\item 反证，假设存在$z_1\ne z_2$，使得$f\left( z_1 \right) =z_1$，$f\left( z_2 \right) =z_2$.令
\begin{align*}
\zeta \left( z \right) =\frac{z_1-z}{1-\overline{z_1}z},\quad g\left( z \right) =\frac{f\left( z_1 \right) -f\left( z \right)}{1-\overline{f\left( z_1 \right) }f\left( z \right)},
\end{align*}
则$\zeta $都是$\mathbb{D}$上的全纯自同构，$g\left( \zeta \left( z \right) \right) \in H\left( \mathbb{D} \right)$.并且
\begin{align*}
g\left( \zeta \left( 0 \right) \right) =0,\quad \left| g\left( \zeta \left( z \right) \right) \right|\leqslant 1,\quad \forall z\in \mathbb{D}.
\end{align*}
由\hyperref[theorem:Schwarz引理-定理4.5.3]{Schwarz引理}知
\begin{align*}
\left| g\left( \zeta \left( z \right) \right) \right|\leqslant \left| z \right|\Longleftrightarrow \left| g\left( z \right) \right|\leqslant \left| \zeta ^{-1}\left( z \right) \right|\Longleftrightarrow \left| \frac{f\left( z_1 \right) -f\left( z \right)}{1-\overline{f\left( z_1 \right) }f\left( z \right)} \right|\leqslant \left| \frac{z_1-z}{1-\overline{z_1}z} \right|.
\end{align*}
由假设知，当$z=z_2$时上式等号成立.故存在$\theta \in \mathbb{R}$，使得对$\forall z\in \mathbb{D}$，有
\begin{align*}
g\left( \zeta \left( z \right) \right) =e^{\mathrm{i}\theta}z\Longleftrightarrow g\left( z \right) =e^{\mathrm{i}\theta}\zeta ^{-1}\left( z \right)\Longleftrightarrow \frac{f\left( z_1 \right) -f\left( z \right)}{1-\overline{f\left( z_1 \right) }f\left( z \right)}=e^{\mathrm{i}\theta}\frac{z_1-z}{1-\overline{z_1}z}.
\end{align*}
将$z=z_2$代入上式得$e^{\mathrm{i}\theta}=1$.于是
\begin{align*}
\zeta \left( f\left( z \right) \right) =\frac{z_1-f\left( z \right)}{1-\overline{z_1}f\left( z \right)}=\frac{f\left( z_1 \right) -f\left( z \right)}{1-\overline{f\left( z_1 \right) }f\left( z \right)}=\frac{z_1-z}{1-\overline{z_1}z}=\zeta \left( z \right) ,\quad \forall z\in \mathbb{D} .
\end{align*}
再由$\zeta \left( z \right)$是$\mathbb{D}$上的全纯自同构知
\begin{align*}
f\left( z \right) =z,\quad \forall z\in \mathbb{D} .
\end{align*}
这与$f$不是恒等映射矛盾!故结论成立.

\item $f\left( z \right) =\frac{z+1}{2};\quad f\left( z \right) =\frac{\frac{1}{2}-z}{1-\frac{1}{2}z}=\frac{1-2z}{2-z}.$
\end{enumerate}
\end{proof}

\begin{example}
设$\mathbb{D} =\left\{ z:\left| z \right|<1 \right\}$,$f:\mathbb{D}\to\mathbb{D}$全纯，满足$f(0)=a,f(a)=0$，其中$a\in\mathbb{D}\setminus\{0\}$. 证明$f$的表达式为
\begin{align*}
f(z)=\frac{a-z}{1-\overline{a}z}.
\end{align*}
\end{example}
\begin{proof}
令
\begin{align*}
g\left( z \right) =\frac{f\left( a \right) -f\left( z \right)}{1-\overline{f\left( a \right) }f\left( z \right)},\quad \zeta \left( z \right) =\frac{a-z}{1-\overline{a}z},
\end{align*}
则$\zeta$是$\mathbb{D}$上的全纯自同构，$g\left( \zeta \left( z \right) \right) \in H\left( \mathbb{D} \right)$，并且
\begin{align*}
g\left( \zeta \left( 0 \right) \right) =0,\quad \left| g\left( \zeta \left( z \right) \right) \right|\leqslant 1,\quad \forall z\in \mathbb{D} .
\end{align*}
由\hyperref[theorem:Schwarz引理-定理4.5.3]{Schwarz引理}知，对$\forall z\in \mathbb{D}$，有
\begin{align*}
\left| g\left( \zeta \left( z \right) \right) \right|\leqslant \left| z \right|\Longleftrightarrow \left| f\left( z \right) \right|=\left| \frac{f\left( a \right) -f\left( z \right)}{1-\overline{f\left( a \right) }f\left( z \right)} \right|=\left| g\left( z \right) \right|\leqslant \left| \zeta ^{-1}\left( z \right) \right|=\left| \zeta \left( z \right) \right|=\left| \frac{a-z}{1-\overline{a}z} \right|.
\end{align*}
当$z=0$时，上述不等式取等，故存在$\theta \in \mathbb{R}$，使得
\begin{align*}
f\left( z \right) =e^{\mathrm{i}\theta}\frac{a-z}{1-\overline{a}z},\quad \forall z\in \mathbb{D} .
\end{align*}
再将$f\left( 0 \right) =a$代入上式得$e^{\mathrm{i}\theta}=1$，故
\begin{align*}
f\left( z \right) =\frac{a-z}{1-\overline{a}z},\quad \forall z\in \mathbb{D} .
\end{align*}
\end{proof}

\begin{example}
设$\mathbb{D} =\left\{ z:\left| z \right|<1 \right\}$,$f:\mathbb{D}\to\mathbb{D}$ 全纯，取定 $a\in\mathbb{D}$，求 $|f'(a)|$ 的最大值，并写出极值映射的一般形式.
\end{example}
\begin{remark}
$f$取极值只需满足等号成立的充分性条件,因此对任意的$\theta\in \mathbb{R}$都成立.
\end{remark}
\begin{proof}
令
\begin{align*}
\varphi_{f(a)}(z)=\frac{f(a)-z}{1-\overline{f(a)}z},\quad \varphi_a(z)=\frac{a-z}{1-\overline{a}z},
\end{align*}
则$\varphi_{f(a)},\varphi_a$都是$\mathbb{D}$上的全纯自同构，且
\begin{align*}
\varphi_{f(a)}(f(a))=\varphi_a(a)=0,\quad \varphi_a'(z)=\frac{a^2-1}{(1-\overline{a}z)^2},\quad \varphi_{f(a)}'(z)=\frac{f^2(a)-1}{(1-\overline{f(a)}z)^2}.
\end{align*}
于是$\varphi_{f(a)}\circ f\circ \varphi_a^{-1}\in H(\mathbb{D})$，且
\begin{align*}
\left|\varphi_{f(a)}\circ f\circ \varphi_a^{-1}(z)\right|\leqslant 1,\quad \varphi_{f(a)}\circ f\circ \varphi_a^{-1}(0)=0.
\end{align*}
由\hyperref[theorem:Schwarz引理-定理4.5.3]{Schwarz引理}知
\begin{align*}
\left|\left(\varphi_{f(a)}\circ f\circ \varphi_a^{-1}\right)'(0)\right|\leqslant 1&\Longleftrightarrow \left|\varphi_{f(a)}'\left[f(\varphi_a(0))\right]\cdot f'(\varphi_a(0))\cdot \varphi_a'(0)\right|=\left|\varphi_{f(a)}'\left[f(a)\right]\cdot f'(a)\cdot \varphi_a'(0)\right|\leqslant 1\\
&\Longleftrightarrow \left|\frac{1}{f^2(a)-1}\cdot f'(a)\cdot (a^2-1)\right|\leqslant 1\Longleftrightarrow \left|f'(a)\right|\leqslant \left|\frac{f^2(a)-1}{a^2-1}\right|=\frac{1-f^2(a)}{1-a^2},
\end{align*}
等号成立当且仅当存在$\theta\in\mathbb{R}$，使得$\varphi_{f(a)}\circ f\circ \varphi_a^{-1}(z)=e^{\mathrm{i}\theta}z$成立。而
\begin{align*}
\frac{1-f^2(a)}{1-a^2}\leqslant \frac{1}{1-a^2},
\end{align*}
等号成立当且仅当$f(a)=0$。因此
\begin{align*}
\max\limits_{a\in\mathbb{D}}\left|f'(a)\right|=\frac{1}{1-a^2}.
\end{align*}
极值映射$f$满足
\begin{align*}
\varphi_{f(a)}\circ f\circ \varphi_a^{-1}(z)=e^{\mathrm{i}\theta}z\text{且}f(a)=0,
\end{align*}
即
\begin{align*}
\varphi_{f(a)}(f(z))=e^{\mathrm{i}\theta}\varphi_a(z)\Longleftrightarrow -f(z)=\frac{f(a)-f(z)}{1-\overline{f(a)}f(z)}=e^{\mathrm{i}\theta}\frac{a-z}{1-\overline{a}z}.
\end{align*}
故极值函数为
\begin{align*}
f(z)=e^{\mathrm{i}\theta}\frac{a-z}{1-\overline{a}z},\quad \theta\text{为任意实数}.
\end{align*}
\end{proof}

\begin{example}
设$\mathbb{D} =\left\{ z:\left| z \right|<1 \right\}$,$f:\mathbb{D}\to\mathbb{D}$ 全纯.
\begin{enumerate}[(1)]
\item 证明: $|f(z)-f(-z)|\leqslant 2|z|$, $z\in\mathbb{D}$.
\item 上式对某 $z_0\in\mathbb{D}\setminus\{0\}$ 成立等号的充要条件是什么?
\end{enumerate}
\end{example}
\begin{proof}
令$g(z)=\frac{f(z)-f(-z)}{2}$，则$g\in H(\mathbb{D})$且$g(0)=0$。由\hyperref[theorem:Schwarz引理-定理4.5.3]{Schwarz引理}知
\begin{align*}
\left|\frac{f(z)-f(-z)}{2}\right|=|g(z)|\leqslant |z|\Longleftrightarrow |f(z)-f(-z)|\leqslant 2|z|,\quad \forall z\in\mathbb{D}.
\end{align*}
上式对某$z_0\in\mathbb{D}\setminus\{0\}$成立等号当且仅当存在$\theta\in\mathbb{R}$，使得$g(z)=e^{\mathrm{i}\theta}z$，也即
\begin{align*}
f(z)-f(-z)=2e^{\mathrm{i}\theta}z.
\end{align*}
\end{proof}

\begin{example}
设$\mathbb{D} =\left\{ z:\left| z \right|<1 \right\}$,$\mathbb{D} \to \mathbb{C}$ 全纯, $f(0) = -1$, 并且满足 $|1 + f(z)| < 1 + |f(z)|$, $\forall z \in \mathbb{D}$. 求 $|f'(0)|$ 的最大值.
\end{example}
\begin{proof}
首先对原不等式 $|1 + f(z)| < 1 + |f(z)|$ 两边平方得到
\begin{align*}
1 + 2\mathrm{Re}f(z) + |f(z)|^2 = |1 + f(z)|^2  < (1 + |f(z)|)^2 = 1 + 2|f(z)| + |f(z)|^2,
\end{align*}
进而
\begin{align*}
\mathrm{Re}f(z) < |f(z)|, \quad \forall z \in \mathbb{D}.
\end{align*}
若$f(z)\in [0,+\infty)$,则$\mathrm{Re}f(z)=|f(z)|$,这与上式矛盾!故$f(\mathbb{D}) \subseteq \mathbb{C} \setminus [0,+\infty).$而 $\mathbb{C} \setminus [0,+\infty)$ 是单连通区域,$f(\mathbb{D}) \subseteq \mathbb{C} \setminus [0,+\infty)$,故$\sqrt{z}$在$f(\mathbb{D})$上可以分出单值解析分支,取一个单值解析分支$g(z)$满足 $g(-1) = \mathrm{i}$,则
\begin{align*}
g\circ f\in H(\mathbb{D} ),\quad g\circ f(0)=\mathrm{i},\quad g\circ f(\mathbb{D} )=g(\mathbb{C} \setminus [0,+\infty ))=\{z:\,\mathrm{Im}z>0\}.
\end{align*}
再令$\varphi(z)=\frac{z-\mathrm{i}}{z+\mathrm{i}}$，则
\begin{align*}
\varphi\circ g\circ f\in H(\mathbb{D}),\quad \varphi\circ g\circ f(0)=0,\quad \varphi\circ g\circ f(\mathbb{D})=\varphi(\{z:\,\mathrm{Im}z>0\})=\mathbb{D}.
\end{align*}
由\hyperref[theorem:Schwarz引理-定理4.5.3]{Schwarz引理}知
\begin{align*}
&\left|(\varphi\circ g\circ f)'(0)\right|\leqslant 1\Longleftrightarrow \left|\varphi'\left[g(f(0))\right]\cdot g'(f(0))\cdot f'(0)\right|\leqslant 1\\
&\Longleftrightarrow \left|\frac{2\mathrm{i}}{(z+\mathrm{i})^2}\bigg|_{z=\mathrm{i}}\cdot \frac{1}{2\sqrt{z}}\bigg|_{z=-1}\cdot f'(0)\right|\leqslant 1\Longleftrightarrow \frac{|f'(0)|}{4}\leqslant 1\Longleftrightarrow |f'(0)|\leqslant 4,
\end{align*}
等号成立当且仅当存在$\theta\in\mathbb{R}$，使得对$\forall z\in\mathbb{D}$，有
\begin{align*}
\varphi\circ g\circ f(z)=e^{\mathrm{i}\theta}z\Longleftrightarrow \frac{\sqrt{f(z)}-\mathrm{i}}{\sqrt{f(z)}+\mathrm{i}}=e^{\mathrm{i}\theta}z\Longleftrightarrow f(z)=-\left(\frac{1+e^{\mathrm{i}\theta}z}{1-e^{\mathrm{i}\theta}z}\right)^2.
\end{align*}
故$\max \left| f' \left( 0 \right) \right|=4$.
\end{proof}

\begin{example}
设$\mathbb{D} =\left\{ z:\left| z \right|<1 \right\}$,$f$ 在 $\mathbb{D}$ 上全纯, $|\mathrm{Re}f(z)| < 1$, $f(0)=0$. 证明
\begin{align*}
|f'(0)| \leqslant \frac{4}{\pi}.
\end{align*}
等号成立的充要条件是什么?
\end{example}
\begin{proof}
由题可知$f(\mathbb{D})=\{z:\,|\mathrm{Re}\,z|<1\}$.令
\begin{align*}
\varphi_1(z)=\frac{\pi}{2}z+\frac{\pi}{2},\quad \varphi_2(z)=\mathrm{i}z,\quad \varphi_3(z)=e^z,\quad \varphi_4(z)=\frac{z-\mathrm{i}}{z+\mathrm{i}},
\end{align*}
则$\varphi_1,\varphi_2,\varphi_3,\varphi_4$都是共形映射,并且
\begin{align*}
&\varphi_4\circ \varphi_3\circ \varphi_2\circ \varphi_1\circ f(0)=0,\quad \varphi_1\circ f(\mathbb{D})=\{z:\,0<\mathrm{Re}\,z<\pi\},\\
&\varphi_2\circ \varphi_1\circ f(\mathbb{D})=\{z:\,0<\mathrm{Im}\,z<\pi\},\\
&\varphi_3\circ \varphi_2\circ \varphi_1\circ f(\mathbb{D})=\{z:\,\mathrm{Im}\,z>0\},\quad \varphi_4\circ \varphi_3\circ \varphi_2\circ \varphi_1\circ f(\mathbb{D})=\{z:|z|<1\}.
\end{align*}
从而由\hyperref[theorem:Schwarz引理-定理4.5.3]{Schwarz引理}知
\begin{align*}
\left|\left(\varphi_4\circ \varphi_3\circ \varphi_2\circ \varphi_1\circ f\right)'(0)\right|\leqslant 1\Longleftrightarrow \left|-\frac{\mathrm{i}}{2}\cdot e^{\frac{\pi}{2}\mathrm{i}}\cdot \mathrm{i}\cdot \frac{\pi}{2}\cdot f'(0)\right|\leqslant 1\Longleftrightarrow |f'(0)|\leqslant \frac{4}{\pi},
\end{align*}
等号成立当且仅当存在$\theta\in\mathbb{R}$,使得
\begin{align*}
\varphi_4\circ \varphi_3\circ \varphi_2\circ \varphi_1\circ f(z)=e^{\mathrm{i}\theta}z,\quad \forall z\in\mathbb{C}.
\end{align*}
也即
\begin{align*}
e^{\mathrm{i}\theta}z=\frac{e^{\mathrm{i}\left(\frac{\pi}{2}f(z)+\frac{\pi}{2}\right)}-\mathrm{i}}{e^{\mathrm{i}\left(\frac{\pi}{2}f(z)+\frac{\pi}{2}\right)}+\mathrm{i}}=\frac{e^{\mathrm{i}\frac{\pi}{2}f(z)}-1}{e^{\mathrm{i}\frac{\pi}{2}f(z)}+1}=\mathrm{i}\tan\frac{\pi f(z)}{4},\quad \forall z\in\mathbb{C}.
\end{align*}
\end{proof}

\begin{example}[上半平面的Schwarz-Pick定理]\,\,
记 $\mathbb{H} = \{z = x + \mathrm{i}y \mid y > 0\}$ 为上半平面，$f:\mathbb{H} \to \mathbb{H}$ 全纯，对任意 $z,w \in \mathbb{H}$，证明
\begin{align*}
\left|\frac{f(z)-f(w)}{f(z)-\overline{f(w)}}\right| \leqslant \left|\frac{z-w}{z-\overline{w}}\right|.
\end{align*}
由此说明
\begin{align*}
|f'(w)| \leqslant \frac{\mathrm{Im}(f(w))}{\mathrm{Im}(w)}.
\end{align*}
\end{example}
\begin{remark}
如何构造$g$：已知分式线性变换$\zeta(z)=\frac{z-w}{z-\overline{w}}$满足$\zeta:\mathbb{H}\to\mathbb{D}$且$\zeta(w)=0$，则$\zeta$的反函数$g$就满足$g:\mathbb{D}\to\mathbb{H}$且$g(0)=w$。
\end{remark}
\begin{proof}
对$\forall w\in\mathbb{H}$，固定$w$。令
\begin{align*}
g(z)=\frac{\overline{w}z-w}{z-1},\quad h(z)=\frac{z-f(w)}{z-\overline{f(w)}},
\end{align*}
则$g,h$都是共形映射。记$\mathbb{D}=\{z:\,|z|<1\}$，则$h\circ f\circ g^{-1}$是$\mathbb{D}\to\mathbb{D}$上的全纯函数，并且
\begin{align*}
h\circ f\circ g(0)=0,\quad h\circ f\circ g(\mathbb{D})=\mathbb{D},\quad g^{-1}(z)=\frac{z-w}{z-\overline{w}}.
\end{align*}
由\hyperref[theorem:Schwarz引理-定理4.5.3]{Schwarz引理}可知
\begin{align*}
|h\circ f\circ g(z)|\leqslant|z|,\quad \forall z\in\mathbb{D} \Longleftrightarrow |h\circ f(z)|\leqslant|g^{-1}(z)|,\quad \forall z\in\mathbb{H}.
\end{align*}
此即
\begin{align*}
\left| \frac{f(z)-f(w)}{f(z)-\overline{f(w)}} \right|\leqslant\left| \frac{z-w}{z-\overline{w}} \right|,\quad \forall z\in\mathbb{H}.
\end{align*}
于是
\begin{align*}
\left| \frac{f(z)-f(w)}{z-w} \right|\leqslant\left| \frac{f(z)-\overline{f(w)}}{z-\overline{w}} \right|,\quad \forall z\in\mathbb{H}.
\end{align*}
由于$f'(w)$存在，故上式两边令$z$沿径向趋于$w$得
\begin{align*}
|f'(w)|\leqslant\left| \frac{f(w)-\overline{f(w)}}{w-\overline{w}} \right|=\frac{\mathrm{Im}f(w)}{\mathrm{Im}w}.
\end{align*}
再由$w$的任意性知结论成立.
\end{proof}

\begin{example}
设$\mathbb{D} =\left\{ z:\left| z \right|<1 \right\}$,$f:\mathbb{D}\to\mathbb{H}_R$ 全纯且 $f(0)=1$，其中 $\mathbb{H}_R=\{z\mid\mathrm{Re}\,z>0\}$ 为右半平面。证明
\begin{align*}
\frac{1-|z|}{1+|z|}\leqslant\mathrm{Re}f(z)\leqslant|f(z)|\leqslant\frac{1+|z|}{1-|z|},\quad \forall z\in\mathbb{D}.
\end{align*}
上式对某个 $z_0\neq0$ 成立的充要条件是什么？
\end{example}
\begin{proof}
令
\begin{align*}
g(z)=\mathrm{i}z,\quad h(z)=\frac{z-\mathrm{i}}{z+\mathrm{i}},
\end{align*}
则$g,h$都是共形映射。从而$h\circ g\circ f$是$\mathbb{D}$上的全纯函数，且
\begin{align*}
h\circ g\circ f(0)=0,\quad h\circ g\circ f(\mathbb{D})=\mathbb{D}.
\end{align*}
由\hyperref[theorem:Schwarz引理-定理4.5.3]{Schwarz引理}知，对$\forall z\in\mathbb{D}$，有
\begin{align*}
\left| h\circ g\circ f(z) \right|\leqslant\left| z \right|\Longleftrightarrow \left| \frac{f(z)-1}{f(z)+1} \right|=\left| \frac{\mathrm{i}f(z)-\mathrm{i}}{\mathrm{i}f(z)+\mathrm{i}} \right|\leqslant\left| z \right|.
\end{align*}
两边同时平方得，对$\forall z\in\mathbb{D}$，有
\begin{align*}
\left| \frac{f(z)-1}{f(z)+1} \right|^2\leqslant\left| z \right|^2&\Longleftrightarrow \left| f(z) \right|^2+1-2\mathrm{Re}f(z) \leqslant \left| z \right|^2\left( \left| f(z) \right|^2+1+2\mathrm{Re}f(z) \right) \\
&\Longleftrightarrow \left( 1-\left| z \right|^2 \right) \left| f(z) \right|^2+1-\left| z \right|^2\leqslant 2\left( \left| z \right|^2+1 \right) \mathrm{Re}f(z).
\end{align*}
再利用$\mathrm{Re}f(z)\leqslant\left| f(z) \right|$，得
\begin{align*}
&\quad \quad \,\,\left( 1-\left| z \right|^2 \right) \left[ \mathrm{Re}f(z) \right] ^2+1-\left| z \right|^2\leqslant \left( 1-\left| z \right|^2 \right) \left| f(z) \right|^2+1-\left| z \right|^2\leqslant 2\left( \left| z \right|^2+1 \right) \mathrm{Re}f(z) \leqslant 2\left( \left| z \right|^2+1 \right) \left| f(z) \right| \\
&\Longleftrightarrow \left( 1-\left| z \right|^2 \right) \left[ \mathrm{Re}f(z) \right] ^2-2\left( \left| z \right|^2+1 \right) \mathrm{Re}f(z)+1-\left| z \right|^2\leqslant 0,\quad \left( 1-\left| z \right|^2 \right) \left| f(z) \right|^2-2\left( \left| z \right|^2+1 \right) \left| f(z) \right|+1-\left| z \right|^2\leqslant 0.
\end{align*}
根据二次函数性质可得
\begin{align}\label{eq::20-jr90wjfjdejg3tgdsgdf}
\frac{1-\left| z \right|}{1+\left| z \right|}\leqslant \mathrm{Re}f(z) \leqslant \left| f(z) \right|\leqslant \frac{1+\left| z \right|}{1-\left| z \right|},\quad \forall z\in\mathbb{D}.
\end{align}
综上可知，上式对某个$z_0\neq 0$等号成立的充要条件是
\begin{align*}
\mathrm{Re}f(z_0)=\left| f(z_0) \right|\text{且}\exists \theta \in \mathbb{R},\,h\circ g\circ f(z)=e^{\mathrm{i}\theta}z\Longleftrightarrow \mathrm{Re}f(z_0)=\left| f(z_0) \right|\text{且}\exists \theta \in \mathbb{R},\,\frac{f(z)-1}{f(z)+1}=e^{\mathrm{i}\theta}z \\
\Longleftrightarrow \mathrm{Re}f(z_0)=\left| f(z_0) \right|\text{且}\exists \theta \in \mathbb{R},\,f(z)=\frac{1+e^{\mathrm{i}\theta}z}{1-e^{\mathrm{i}\theta}z}\Longleftrightarrow \exists \theta \in \mathbb{R},\,a\geqslant 0,\,f(z)=\frac{1+e^{\mathrm{i}\theta}z}{1-e^{\mathrm{i}\theta}z}\text{且}f(z_0)=\frac{1+e^{\mathrm{i}\theta}z_0}{1-e^{\mathrm{i}\theta}z_0}=a.
\end{align*}
注意到
\begin{align*}
\frac{1+e^{\mathrm{i}\theta}z_0}{1-e^{\mathrm{i}\theta}z_0}=a\Longleftrightarrow e^{\mathrm{i}\theta}=\frac{a-1}{(a+1)z_0},
\end{align*}
两边同时取模得$1=\frac{a-1}{(a+1)\left| z_0 \right|}$，进而$a=\frac{\left| z_0 \right|+1}{1-\left| z_0 \right|}$，于是$e^{\mathrm{i}\theta}=\frac{a-1}{(a+1)z_0}=\frac{\left| z_0 \right|}{z_0}$。故\eqref{eq::20-jr90wjfjdejg3tgdsgdf}式对某个$z_0\neq 0$等号成立的充要条件是
\begin{align*}
f(z)=\frac{1+\frac{\left| z_0 \right|}{z_0}z}{1-\frac{\left| z_0 \right|}{z_0}z}=\frac{z_0+\left| z_0 \right|z}{z_0-\left| z_0 \right|z},\quad \forall z\in\mathbb{D}.
\end{align*}
\end{proof}

\begin{theorem}[Schwarz-Pick定理的一般形式]\label{theorem:Schwarz-Pick定理的一般形式}
设$\mathbb{D}=\{z:\,|z|<1\}$,$f: \mathbb{D} \to \mathbb{D}$ 全纯, $a$ 是 $f(z)-f(a)$ 的 $n(n\geqslant 1)$ 阶零点,则
\begin{enumerate}[(1)]
\item 对任意 $z\in\mathbb{D}$, 成立
\begin{align*}
\left|\frac{f(a)-f(z)}{1-\overline{f(a)}f(z)}\right|\leqslant\left|\frac{a-z}{1-\overline{a}z}\right|^n.
\end{align*}
存在$z_0\in\mathbb{D},z_0\neq 0$，使得$\left|\frac{f(z_0)-f(a)}{1-\overline{f(a)}f(z_0)}\right|=\left|\frac{z_0-a}{1-\overline{a}z_0}\right|^n$,的充要条件是存在$\theta\in[0,2\pi]$,使得
\begin{align*}
\frac{f(a)-f(z)}{1-\overline{f(a)}f(z)}=e^{\mathrm{i}\theta}\left(\frac{a-z}{1-\overline{a}z}\right)^n,\quad \forall z\in\mathbb{D}.
\end{align*}
并且此时$f\in\mathrm{Aut}(\mathbb{D})$.

\item $f^{(n)}(a)$ 满足如下估计
\begin{align*}
|f^{(n)}(a)|\leqslant n!\frac{1-|f(a)|^2}{(1-|a|^2)^n}.
\end{align*}
$|f^{(n)}(a)|=n!\frac{1-|f(a)|^2}{(1-|a|^2)^n}$成立的充要条件是存在$\theta\in[0,2\pi]$,使得
\begin{align*}
\frac{f(a)-f(z)}{1-\overline{f(a)}f(z)}=e^{\mathrm{i}\theta}\left(\frac{a-z}{1-\overline{a}z}\right)^n,\quad \forall z\in\mathbb{D}.
\end{align*}
并且此时$f\in\mathrm{Aut}(\mathbb{D})$.
\end{enumerate}
\end{theorem}
\begin{proof}
\begin{enumerate}[(1)]
\item 令
\begin{align*}
\varphi(z)=\frac{a-z}{1-\overline{a}z},\quad \psi(z)=\frac{f(a)-z}{1-\overline{f(a)}z},\quad g(z)=\psi\circ f\circ \varphi(z),
\end{align*}
则$\varphi,\psi\in\mathrm{Aut}(\mathbb{D})$，$g\in H(\mathbb{D})$且$g(0)=0$，$g(\mathbb{D})\subseteq\mathbb{D}$。由条件知，存在$f_1\in H(\mathbb{D})$，使得
\begin{align*}
f(z)=f(a)+(z-a)^n f_1(z).
\end{align*}
从而
\begin{align*}
g(z)&=\psi\circ f\circ \varphi(z)=\psi\left[f(a)+\left(\varphi(z)-a\right)^n f_1\left(\varphi(z)\right)\right]\\
&=\psi\left[f(a)+\frac{z^n(a^2-1)^n}{1-\overline{a}z}f_1\left(\varphi(z)\right)\right]=\frac{\frac{z^n(a^2-1)^n}{1-\overline{a}z}f_1\left(\varphi(z)\right)}{1-\overline{f(a)}f\left(\varphi(z)\right)}\\
&=z^n\frac{(a^2-1)^n f_1\left(\varphi(z)\right)}{(1-\overline{a}z)\left[1-\overline{f(a)}f\left(\varphi(z)\right)\right]}.
\end{align*}
而$h(z)=\frac{(a^2-1)^n f_1\left(\varphi(z)\right)}{(1-\overline{a}z)\left[1-\overline{f(a)}f\left(\varphi(z)\right)\right]}\in H(\mathbb{D})$，因此$z=0$是$g$的$n$阶零点。也即
\begin{align*}
h(z)=\frac{g(z)}{z^n}\in H(\mathbb{D}).
\end{align*}
对$\forall z\in\mathbb{D}$，再对$\forall r\in(|z|,1)$，由\hyperref[theorem:最大模原理-定理4.5.1]{最大模原理}得
\begin{align*}
|h(z)|\leqslant\max\limits_{|z|=r}\frac{|g(z)|}{r^n}\leqslant\frac{1}{r^n}.
\end{align*}
令$r\rightarrow 1^-$，得$|h(z)|\leqslant 1,\forall z\in\mathbb{D}$。故$|g(z)|\leqslant |z|^n,\forall z\in\mathbb{D}$。即对$\forall z\in\mathbb{D}$，有
\begin{align*}
|\psi\circ f\circ \varphi(z)|\leqslant |z|^n\Longleftrightarrow |\psi\circ f(z)|\leqslant |\varphi^{-1}(z)|^n=|\varphi(z)|^n\Longleftrightarrow \left|\frac{f(a)-f(z)}{1-\overline{f(a)}f(z)}\right|\leqslant \left|\frac{a-z}{1-\overline{a}z}\right|^n.
\end{align*}
若存在$z_0\in\mathbb{D},z_0\neq 0$，使得$\left|\frac{f(z_0)-f(a)}{1-\overline{f(a)}f(z_0)}\right|=\left|\frac{z_0-a}{1-\overline{a}z_0}\right|^n$，即$|h(z_0)|=1$。这说明全纯函数$h$在内点$z_0$处取到了最大模$1$，由\hyperref[theorem:最大模原理-定理4.5.1]{最大模原理}知，$h$必是常数。设$h(z)=c\in\mathbb{C}$，由$|h(z_0)|=1$得$|c|=1$，故存在$\theta\in[0,2\pi]$，使得$c=e^{\mathrm{i}\theta}$，因而对$\forall z\in\mathbb{D}$，有
\begin{align*}
g(z)&=z^n h(z)=e^{\mathrm{i}\theta}z^n\Longleftrightarrow \psi\circ f\circ \varphi(z)=e^{\mathrm{i}\theta}z^n\\
&\Longleftrightarrow \psi\circ f(z)=e^{\mathrm{i}\theta}\left(\varphi^{-1}(z)\right)^n=e^{\mathrm{i}\theta}\varphi^n(z)\\
&\Longleftrightarrow \frac{f(a)-f(z)}{1-\overline{f(a)}f(z)}=e^{\mathrm{i}\theta}\left(\frac{a-z}{1-\overline{a}z}\right)^n.
\end{align*}
此时$f\in\mathrm{Aut}(\mathbb{D})$。

\item  由(1)可得
\begin{align*}
\left|\frac{f(z)-f(a)}{1-\overline{f(a)}f(z)}\right|\leqslant \left|\frac{z-a}{1-\overline{a}z}\right|^n,\quad \forall z\in\mathbb{D}.
\end{align*}
从而
\begin{align*}
\left|\frac{f(z)-f(a)}{(z-a)^n}\right|\leqslant \left|\frac{1-\overline{f(a)}f(z)}{(1-\overline{a}z)^n}\right|,\quad \forall z\in\mathbb{D}.
\end{align*}
令$z\rightarrow a$，由\hyperref[theorem:解析函数一般形式的L'Hospital法则]{L'Hospital法则}可得
\begin{align*}
\left|\frac{f^{(n)}(a)}{n!}\right|=\lim\limits_{z\rightarrow a}\left|\frac{f(z)-f(a)}{(z-a)^n}\right|\leqslant\lim\limits_{z\rightarrow a}\left|\frac{1-\overline{f(a)}f(z)}{(1-\overline{a}z)^n}\right|=\frac{1-|f(a)|^2}{(1-|a|^2)^n}.
\end{align*}
故
\begin{align*}
|f^{(n)}(a)|\leqslant n!\frac{1-|f(a)|^2}{(1-|a|^2)^n}.
\end{align*}
若$|f^{(n)}(a)|=n!\frac{1-|f(a)|^2}{(1-|a|^2)^n}$，令
\begin{align*}
H(z)=\frac{f(z)-f(a)}{(z-a)^n}\cdot\frac{(1-\overline{a}z)^n}{1-\overline{f(a)}f(z)},\quad \forall z\in\mathbb{D},
\end{align*}
则由上述证明知$|H(z)|\leqslant 1$，又由假设知$|H(a)|=\lim\limits_{z\rightarrow a}|H(z)|=\frac{|f^{(n)}(a)|}{n!}\cdot\frac{(1-|a|^2)^n}{1-|f(a)|^2}=1$。而$a$是$\mathbb{D}$的内点，由\hyperref[theorem:最大模原理-定理4.5.1]{最大模原理}知，$H$必为常数，设$H(z)=C\in\mathbb{C}$，则由$|H(a)|=1$知$|C|=1$，因此存在$\theta\in[0,2\pi]$，使得$H(z)=e^{\mathrm{i}\theta}$，即
\begin{align*}
\frac{f(a)-f(z)}{1-\overline{f(a)}f(z)}=e^{\mathrm{i}\theta}\left(\frac{a-z}{1-\overline{a}z}\right)^n,\quad \forall z\in\mathbb{D}.
\end{align*}
此时$f\in\mathrm{Aut}(\mathbb{D})$。
\end{enumerate}
\end{proof}

\begin{theorem}[Schwarz定理的一般形式]\label{theorem:Schwarz定理的一般形式}
设$\mathbb{D} = \{z: |z| < 1\}$,$f: \mathbb{D} \to \mathbb{D}$全纯，且$z=0$是$f(z)$的$n$阶零点，则
\begin{enumerate}[(1)]
\item 对任意$z\in\mathbb{D}$，有
\begin{align*}
|f(z)| \leqslant |z|^n.
\end{align*}
存在$z_0\in\mathbb{D}, z_0\neq 0$使得$|f(z_0)|=|z_0|^n$的充要条件是存在$\theta\in[0,2\pi]$，使得
\begin{align*}
f(z) = e^{\mathrm{i}\theta} z^n,\quad \forall z\in\mathbb{D}.
\end{align*}

\item $f$在$0$处的$n$阶导数满足
\begin{align*}
|f^{(n)}(0)| \leqslant n!\,\,.
\end{align*}
$|f^{(n)}(0)|=n!$成立的充要条件是存在$\theta\in[0,2\pi]$，使得
\begin{align*}
f(z) = e^{\mathrm{i}\theta} z^n,\quad \forall z\in\mathbb{D}.
\end{align*}
\end{enumerate}
\end{theorem}
\begin{proof}
在\hyperref[theorem:Schwarz-Pick定理的一般形式]{Schwarz-Pick定理的一般形式}中取$a=f(a)=0$即得.
\end{proof}

\begin{example}
设$f(z)$在$D: |z|<1$内全纯，且$|f(z)|\leqslant 1$（$z\in D$）。若$z_0$是$D$中一点，使得$z_0$和$-z_0$都是$f(z)$的$m$阶零点，且$0<|z_0|\leqslant \frac{m-1}{m}$，则$|f(0)|<e^{-2}$。
\end{example}
\begin{note}
证明思路类似\hyperref[theorem:Schwarz-Pick定理的一般形式]{Schwarz-Pick定理的一般形式}的证明.
直接利用\hyperref[theorem:Schwarz定理的一般形式]{Schwarz定理的一般形式}解决不太方便.
\end{note}
\begin{proof}
令
\begin{align*}
\varphi(z)=\left(\frac{z_0-z}{1-\overline{z_0}z}\cdot\frac{z+z_0}{1+\overline{z_0}z}\right),\quad g(z)=\frac{f(z)}{\varphi^m(z)},
\end{align*}
则由分式线性变换的性质知,$\varphi\in H(D)$且
\begin{align*}
|\varphi(z)|<1,\ \forall z\in D;\quad |\varphi(z)|=1,\ \forall z\in\{z:|z|=1\}.
\end{align*}
由$z_0,-z_0$都是$f(z)$的$m$阶零点知,$g\in H(D)$.对$\forall z\in D$,任取$r\in(|z|,1)$,则由最大模原理知
\begin{align*}
|g(z)|\leqslant\max\limits_{|z|=r}|g(z)|=\max\limits_{|z|=r}\frac{|f(z)|}{|\varphi(z)|}\leqslant\max\limits_{|z|=r}\frac{1}{|\varphi(z)|}.
\end{align*}
令$r\rightarrow1^-$,得
\begin{align*}
|g(z)|\leqslant\lim\limits_{r\rightarrow1^-}\max\limits_{|z|=r}\frac{1}{|\varphi(z)|}=\varlimsup\limits_{|z|\rightarrow1^-}\frac{1}{|\varphi(z)|}=1,\quad \forall z\in D.
\end{align*}
故
\begin{align*}
\left| g\left( 0 \right) \right|=\frac{\left| f\left( 0 \right) \right|}{\left| \varphi ^m\left( 0 \right) \right|}\leqslant 1\Longrightarrow \left| f\left( 0 \right) \right|\leqslant \left| z_0 \right|^{2m}\leqslant \left( \frac{m-1}{m} \right) ^{2m}=\left( 1-\frac{1}{m} \right) ^{2m}<e^{-2}.
\end{align*}
\end{proof}

\begin{theorem}[Borel-Carathéodory定理]\label{theorem:Borel-Carathéodory 定理}
若 $f$ 在 $B(0,R)$ 上全纯, 记
\begin{align*}
M(r) = \max\limits_{|z|=r} |f(z)|, \quad A(r) = \max\limits_{|z|=r} \mathrm{Re}f(z), \quad r \in [0,R].
\end{align*}
\begin{enumerate}[(1)]
\item 如果 $f(0) = 0$, 证明当 $|z| < r$ 时,
\begin{align*}
-\frac{2A(r)|z|}{r - |z|} \leqslant \mathrm{Re}f(z) \leqslant \frac{2A(r)|z|}{r + |z|}; \quad |f(z)| \leqslant \frac{2A(r)|z|}{r - |z|}.
\end{align*}

\item 证明: 对任意 $r \in [0,R)$,
\begin{align*}
M(r) \leqslant \frac{2r}{R - r}A(R) + \frac{R + r}{R - r}|f(0)|, \\
A(r) \leqslant \frac{2r}{R + r}A(R) + \frac{R - r}{R + r}\mathrm{Re}f(0).
\end{align*}
\item 如果 $f$ 不取零值, 证明
\begin{align*}
M(r) \leqslant M(0)^{\frac{R-r}{R+r}} M(R)^{\frac{2r}{R+r}}.
\end{align*}
\end{enumerate}
\end{theorem}
\begin{proof}
\end{proof}

\begin{example}
记$\mathscr{F}$是全纯函数$f:\mathbb{D} \to \mathbb{D}$全体. 对任意$a,b \in \mathbb{D}$, 求极值
\begin{align*}
\rho(a,b) = \sup\limits_{f \in \mathscr{F}} |f(a) - f(b)|.
\end{align*}
\end{example}
\begin{proof}

\end{proof}










































































































































\end{document}